\chapter{投影、插值与 Galerkin 框架}\label{ch:6}
\begin{chaptersummary}
弱形式规定了解与测试函数所在的连续空间，离散计算则进一步选择有限维表示及残差约束。本章据此比较插值、正交投影、Galerkin、配置点与控制体守恒，再用误差正交性联系空间逼近与数值解。后续各方法将在这一共同框架下定义自由度、算子和界面耦合。
\end{chaptersummary}
\learninggoals{写出试探与测试空间；区分五类逼近或残差机制；说明能量投影的适用条件；由局部基函数构造质量矩阵、残差和稀疏耦合。}
\studyroute{这里给出强制问题的基本误差估计，略去一般 Petrov--Galerkin 的离散 inf--sup 理论和非协调误差估计。}{Galerkin orthogonality; Riesz representation; Cea lemma; Petrov--Galerkin inf-sup}{对 $-\!u''=f$ 比较两种线性有限元网格的能量误差，再以对流问题解释测试空间改变后的稳定性问题。}

\section{从连续算子到残差约束}
设连续问题为 $\mathcal A u=f$，允许的场属于 $V$，测试空间为 $W$。
弱形式可写为 $b(u,v)=\ell(v)$。
选择有限维试探空间 $V_h$ 和测试空间 $W_h$，得到
\begin{equation}\label{eq:galerkin-general}
\text{求 }u_h\in V_h:\qquad b(u_h,v_h)=\ell(v_h)
\quad\forall v_h\in W_h.
\end{equation}
残差泛函 $r_h(v)=\ell(v)-b(u_h,v)$ 在 $W_h$ 上为零；
它在其余测试方向上通常仍有误差。相同空间称为 Galerkin；
不同试探和测试空间构成 Petrov--Galerkin。对非线性问题，残差与解的关系相应成为非线性。

\conceptbridge{有限维表示确定之后，还缺少哪一步才能得到数值方程？}
{基函数只给出了候选解的表示。数值方程还需规定哪些测试方向、网格点或控制体上的残差为零。这一约束机制与界面处理共同定义离散算子。}
{表示空间 $\longrightarrow$ 残差约束 $\longrightarrow$ 有限维方程。}

\section{插值与正交投影}\index{Galerkin 方法}\index{插值}
插值算子 $I_h$ 通过指定自由度 $\lambda_i$ 满足
$\lambda_i(I_hu)=\lambda_i(u)$。节点插值要求点值有意义；
在一般高维 $H^1$ 空间中，点值未必是连续泛函，可改用局部平均型自由度。

Hilbert 空间内的正交投影 $P_h$ 则满足
\[
P_hu\in V_h,\qquad (u-P_hu,v_h)_V=0\quad\forall v_h\in V_h.
\]
对 $u(x)=x^2$、$V_h=\Poly_1(-1,1)$，端点线性插值得到 $I_hu=1$，
而 $L^2$ 正交投影得到 $P_hu=1/3$。
同一个表示空间可以对应不同的逼近算子。

\section{Riesz 表示、能量范数与误差估计}
Hilbert 空间的 Riesz 映射 $R_V:V\to V'$ 由
$\langle R_Vw,v\rangle=(w,v)_V$ 定义，
使每个连续线性泛函对应唯一的内积表示向量。具体 PDE 的双线性形式还需核对对称性和强制性。
对 $V_h\subset V$ 的协调 Galerkin 近似，连续与离散方程相减得
\[
b(u-u_h,v_h)=0\quad\forall v_h\in V_h.
\]
若 $b$ 对称正定，它定义能量范数 $\norm{v}_b^2=b(v,v)$，并有
\[
\norm{u-u_h}_b=\inf_{w_h\in V_h}\norm{u-w_h}_b.
\]
对有界且强制、但未必对称的 $b$，由任意 $w_h\in V_h$ 可得
\[
\alpha\norm{u-u_h}_V^2
\le b(u-u_h,u-w_h)
\le M\norm{u-u_h}_V\norm{u-w_h}_V.
\]
约去误差范数即得到 C\'ea 估计
\[
\norm{u-u_h}_V\le\frac{M}{\alpha}
\inf_{w_h\in V_h}\norm{u-w_h}_V.
\]
这一结论把数值误差联系到空间逼近能力。Petrov--Galerkin 与非强制问题需离散 inf--sup 条件。

\section{不同方法的自由度与残差机制}
\begin{center}\small
\begin{tabularx}{\textwidth}{@{}l>{\raggedright\arraybackslash}X>{\raggedright\arraybackslash}X@{}}
\toprule
方法 & 表示与自由度 & 方程或残差约束\\
\midrule
FDM & 网格函数、节点值 & 差分算子在网格点满足方程\\
FVM & 单元平均与局部重构 & 每个控制体的积分守恒\\
FEM & 局部基函数、迹相容自由度 & 对有限维测试空间满足弱形式\\
谱 Galerkin & 全局模式系数 & 对所选全局测试模式残差为零\\
谱配置 & 全局插值多项式的节点值 & 配置点满足微分方程\\
DG / FR & 单元局部多项式 & 弱残差 / 修正通量导数及共享界面通量\\
LBM & 各离散速度的粒子群密度 & 迁移与碰撞及其守恒矩\\
\bottomrule
\end{tabularx}
\end{center}
这些方法可共享有限维表示、残差和结构分析的语言。
Hilbert 正交投影的具体解释由内积、试探空间与残差条件共同决定。

\aiexercise{在 $[-1,1]$ 上分别求 $x^2$ 到 $\Poly_1$ 的端点插值与 $L^2$ 投影，显式计算两个误差。再说明对称 Poisson Galerkin 的正交性使用哪一个内积。}
{比较中心差分、单元平均 FVM 和线性 FEM 对同一 Poisson 问题的自由度及残差约束；列出均匀网格上产生相同矩阵所需的积分和边界条件。}

\section{有限维代数系统的形成}
偏微分方程离散后通常得到
\[
\bm A\bm x=\bm b
\]
或非线性残差系统
\[
\bm R(\bm U)=0.
\]

\begin{authornote}
\textbf{潜在疑问：}
“为什么一个连续 PDE 离散后会变成有限维向量方程？自由度、残差和矩阵条目分别来自哪里？”

\textbf{作者补充：}
连续场在每一点都有取值，计算机只能保存有限个参数。离散化先用网格值或基函数系数表示近似场，再通过节点配点、控制体守恒或对测试函数积分等方式，要求方程在有限个意义下成立。于是连续问题转化为有限维代数方程；非线性或含时问题还需进一步线性化或时间离散。

\textbf{核心概念联系：}
\[
\text{continuous PDE}
\longrightarrow
\text{finite-dimensional DOFs}
\longrightarrow
\text{algebraic system}
\]
\end{authornote}

局部微分算子通常产生稀疏矩阵，边界积分算子常产生稠密矩阵。矩阵结构直接反映物理耦合范围。

\begin{authornote}
\textbf{潜在疑问：}
“为什么 differential operator 的局部性会变成 matrix sparsity？为什么积分方程又会变稠密？”

\textbf{作者补充：}
矩阵的一行记录一个离散方程依赖哪些自由度。局部微分算子配合局部基函数或有限差分模板时，该方程通常只涉及附近自由度，因此每行只有少量非零项。若采用全局基函数或积分算子，一个自由度可能与许多远处自由度耦合，矩阵就可能变稠密；稀疏性取决于算子与离散方式的共同选择。

\textbf{核心概念联系：}
\[
\text{operator locality}
\longrightarrow
\text{matrix sparsity or density}
\]
\end{authornote}


取基函数 $\phi_j$，$u_h=\sum_jU_j\phi_j$，式~\eqref{eq:galerkin-general}~产生
$A_{ij}=b(\phi_j,\psi_i)$ 与 $b_i=\ell(\psi_i)$。
对含时问题，
\[
M_{ij}=(\phi_j,\psi_i),\qquad M_h\dot U_h=R_h(U_h,t).
\]
质量矩阵表示时间导数的配对，残差包含体积、边界及源项。
它们的结构由表示与约束产生；代数求解算法在第 0 卷展开。
